\chapter{MCMC Family II: Population and Differential Evolution}
\label{ch:mcmc-2}

\newthought{These samplers run many chains at once} and build their proposals from the
\emph{population} of chains rather than a fixed step. That makes them good at correlated and
multi-modal posteriors, where a single-chain random walk gets stuck. They share the standard
\textsc{mcmc} \yamlkey{Bounds} from Chapter~\ref{ch:mcmc-1}; the tempering methods add two keys.

\section{Differential-evolution chains}
\label{sec:mcmc2-de}

These three use \yamlkey{num\_chains}, \yamlkey{num\_iters}, and \yamlkey{proposal\_scale}.

\paragraph{DEMCMC} \emph{Differential-evolution \textsc{mcmc}.} It runs a population of chains and
proposes each chain's next step from the \emph{difference} between two other randomly chosen
chains' current states (scaled by \yamlkey{proposal\_scale}), then accepts with the usual
Metropolis rule. Because the proposal direction and length come from the population itself, the
chains automatically adapt to the posterior's shape and correlations.

\paragraph{DREAM} \emph{Differential-evolution adaptive Metropolis}\cite{Vrugt:2009}. A more
powerful population method: it forms proposals from scaled differences of several chain pairs,
adapts a crossover pattern so it can update subsets of parameters at a time, and detects and
resets stuck (outlier) chains. This combination explores high-dimensional, multi-modal posteriors
robustly.

\paragraph{DREAMLite} A lighter-weight \sampler{DREAM} variant with the same configuration and the
same core idea, trading some of the full method's machinery for lower overhead. Also selectable as
\alias{DREAM-lite}{DREAMLite}.

\begin{yamlwin}
Sampling:
  Method: "DREAM"
  Variables: [ ... ]
  Bounds:
    num_chains: 16
    num_iters: 2000
    proposal_scale: 0.5
  LogLikelihood:
    - {name: "LogL", expression: "..."}
\end{yamlwin}

\section{Ensemble sampling}
\label{sec:mcmc2-ensemble}

\paragraph{EnsembleMCMC} An affine-invariant ensemble sampler. It evolves a population of
``walkers''; to move one walker it picks another at random and proposes a point along the line
joining them (a \emph{stretch} move). The trick is that this scheme is invariant to linear
rescalings of the parameters, so it copes with strongly correlated or badly-scaled posteriors with
almost no tuning. Uses the standard \yamlkey{Bounds} (here \yamlkey{num\_chains} is the number of
walkers).

\paragraph{When to use} Correlated posteriors where you would rather not tune a proposal covariance
by hand; a strong general-purpose default.

\section{Parallel tempering}
\label{sec:mcmc2-pt}

When a posterior has well-separated peaks, a single chain rarely crosses the low-likelihood valleys
between them. \emph{Parallel tempering} fixes this by running several chains at different
\emph{temperatures}: hot chains see a flattened likelihood and roam freely across the whole space,
cold chains sample the true posterior, and the chains periodically \emph{swap} states so discoveries
made by the hot chains reach the cold ones.

These methods add two keys to the standard \yamlkey{Bounds}:

\begin{description}[style=nextline, leftmargin=1.2em]
  \item[\yamlkey{exchange\_interval}] how often (in iterations) to attempt swaps between
        neighbouring temperatures.
  \item[\yamlkey{proposal\_scales}] a \emph{list} of proposal scales, one per chain/temperature
        (replacing the single \yamlkey{proposal\_scale}).
\end{description}

\paragraph{PTMCMC} Parallel-tempered random-walk \textsc{mcmc}: one tempered chain per rung of a
temperature ladder, with periodic swaps.

\paragraph{PTEnsemble} The same tempering idea applied to \emph{ensemble} sampling: a tempered
ensemble at each rung, combining affine-invariant moves with mode-hopping swaps.

\begin{yamlwin}
Sampling:
  Method: "PTMCMC"
  Variables: [ ... ]
  Bounds:
    num_chains: 8
    num_iters: 4000
    exchange_interval: 20
    proposal_scales: [0.05, 0.1, 0.2, 0.4, 0.8, 1.2, 1.6, 2.0]
  LogLikelihood:
    - {name: "LogL", expression: "..."}
\end{yamlwin}

\section{Summary}
\label{sec:mcmc2-summary}

For correlated posteriors reach for \sampler{EnsembleMCMC} or \sampler{DREAM}; for multi-modal ones
use the tempered \sampler{PTMCMC}/\sampler{PTEnsemble}. All take more compute than single-chain
\textsc{mcmc} but explore far more reliably.
